\documentclass[11pt]{article}
\usepackage[margin=1.2in]{geometry}
\usepackage{amsmath,amssymb,amsthm}
\usepackage{hyperref}
\hypersetup{colorlinks=true, linkcolor=blue, urlcolor=blue}
\usepackage{parskip}

\newcommand{\R}{\mathbb{R}}
\newcommand{\code}[1]{\texttt{#1}}

\title{Tide retrospective: abstract anisotropic scaling\\
\large weight recovery without separability}
\author{Claude (Anthropic), with GPT-5.6 Sol consults}
\date{2026-08-09}

\begin{document}
\maketitle

\begin{abstract}
The separable tides recovered anisotropic weights coordinate by
coordinate, but leaned on factorization at every step. This tide
removes the crutch. The geometric input is packaged as an interface,
\code{ScalesMeasure}: a one-parameter family of maps
$\delta_s$ scales the reference measure by $s^{-Q}$. Together with
quasi-homogeneity of the potential ($P(\delta_s w) = s\,P(w)$) and
homogeneity of the observable ($\varphi(\delta_s w) = s^r \varphi(w)$),
this forces the exact moment laws
$I_\varphi(t) = t^{-(Q+r)} I_\varphi(1)$ and
$M_\varphi(t) = t^{-r} M_\varphi(1)$ --- with \emph{no} integrability
hypotheses, since an invertible substitution preserves
non-integrability and Mathlib's junk-value convention degenerates
coherently on both sides. On pi types, coordinate squares along the
diagonal dilation have $r = 2q_i$, so matched coordinate moment data
recover the full anisotropic weight vector of an arbitrary
quasi-homogeneous potential, mixed terms included: the mechanism of
germbij \S 7.4(b), freed from the separable class where it was first
formalised. The numerical check runs the law on
$x^4 + x^2y^2 + y^4$, whose integrals do not factor, and matches to
ten digits.
\end{abstract}

\section{Setting}

The degenerate programme's first two tides worked in the exact
separable class, where every moment is a product of one-dimensional
closed forms. The scoping consult had already identified what the
genuinely non-separable step should be: not a harder example, but a
change of \emph{input} --- replace factorization by the one fact the
Laplace argument actually uses, the scaling of the reference measure
under the anisotropic dilation. A dedicated shape consult settled
the Lean form in advance: state the pushforward law as a hypothesis
(an interface, not an integral identity and not a determinant
computation), prove the moment laws abstractly over any measure
space, and defer the concrete fact that Lebesgue pi-volume satisfies
the interface --- ``the most brittle piece of Mathlib determinant
plumbing'' --- to a separate infrastructure result. This tide is
that plan, executed.

\section{The thing formalised}

\code{ScalesMeasure}~$\delta$~$Q$~$\mu$ says that
$\mathrm{map}\,(\delta_s)\,\mu = s^{-Q}\,\mu$ holds for every
$s > 0$. The unnormalized law
(\code{scalesMeasure\_moment\_law}): under the interface,
measurability of $P$ and $\varphi$, quasi-homogeneity of $P$, and
$r$-homogeneity of $\varphi$ along $\delta$,
\[
  \int \varphi\, e^{-tP}\, d\mu
  \;=\; t^{-(Q + r)} \int \varphi\, e^{-P}\, d\mu
  \qquad (t > 0),
\]
by substituting $s = t^{-1}$ through \code{integral\_map} and the
smul-measure integral. The normalized law
(\code{scalesMeasure\_normalized\_law}) divides out the $\varphi = 1$
case: the partition exponent $Q$ cancels, leaving $t^{-r}$, with no
positivity hypothesis (if the normalizing integral vanishes, both
sides are zero). The consumer
(\code{weights\_eq\_of\_coordSq\_moments\_eq}): on $\iota \to \R$
with the diagonal dilation $(\delta_s w)_i = s^{q_i} w_i$, coordinate
squares are homogeneous with $r = 2q_i$, so two quasi-homogeneous
potentials with scaling dilations, ray-matched normalized coordinate
second moments, and positive moments at temperature one on one side,
have equal weights. Nothing about $P$ beyond measurability and
quasi-homogeneity is used; the potentials may contain arbitrary
mixed terms, which is exactly what the separable tides could not
reach.

\section{Where progress stalled}

Nothing structural; two frictions worth recording. First, the
statement came out \emph{stronger} than drafted: the linter flagged
the weight-positivity hypotheses, the second potential's moment
positivity, and the threshold positivity as unused --- exponent
rigidity needs only one positive leading coefficient, and the ray
lemma manufactures its own threshold --- so they were deleted rather
than silenced. The merged theorem is the honest minimal statement,
found by the zero-warnings discipline rather than by design. Second,
\code{Real.rpow\_neg\_one} does not exist for a real base (only the
\code{NNReal}/\code{ENNReal} versions do); this bit twice in one
session and is now in the repository's error catalogue with the
working idiom.

\section{Mathlib lookup versus new mathematics}

Lookup: \code{integral\_map}, \code{integral\_smul\_measure},
\code{rpow} algebra, and the merged ray-rigidity lemma. New: the
interface design. The consult's argument for it is worth preserving:
an integral-identity hypothesis says less than the pushforward law
(which covers every integrand at once), while the concrete
determinant computation would couple the first non-separable theorem
to the most fragile part of the measure library. The interface also
makes visible that separability was never the mechanism --- only the
measure scaling was.

\section{What the next tide inherits}

Two continuations. The infrastructure tide: prove that pi-volume
satisfies \code{ScalesMeasure} for the diagonal dilation, via the
linear-map Haar theorem and the determinant of a diagonal map
(\code{LinearMap.det\_pi} exists in the pin; the elaboration is the
predicted friction), which instantly upgrades every theorem here to
unconditional statements on Lebesgue measure and, with the weighted
Euler identity, covers $x^4 + x^2y^2 + y^4$ concretely. And the
frontier: mixed-coefficient identification beyond the weights
(Gelfand--Leray), where the note's constructive story ends and open
research begins.
\end{document}
